\endinput
%------------------------------------------------------------------------------%
\begin{question}
\end{question}

%------------------------------------------------------------------------------%
\begin{resolution}{Solução}
\end{resolution}

%------------------------------------------------------------------------------%
\endinput


\begin{frame}{Exemplo -- Sistema homogêneo}
Considere
\[
\begin{cases}
x+y+z=0\\
2x+2y+2z=0\\
x-y+z=0.
\end{cases}
\]
\medskip

A matriz dos coeficientes é
\[
A=
\begin{pmatrix}
1&1&1\\
2&2&2\\
1&-1&1
\end{pmatrix}.
\]
\medskip

A segunda linha é o dobro da primeira. Portanto,
\[
\det(A)=0.
\]
\medskip

Logo, existem soluções não triviais.
\end{frame}

\begin{frame}{Exemplo -- Sistema homogêneo}
As equações independentes são
\[
x+y+z=0
\]
e
\[
x-y+z=0.
\]
\medskip

Subtraindo:
\[
2y=0,
\]
portanto,
\[
y=0.
\]
\medskip

Da primeira equação:
\[
x+z=0.
\]
Logo,
\[
x=-z.
\]
\medskip

Tomando
\[
z=t,
\]
temos
\[
(x,y,z)=(-t,0,t).
\]
\end{frame}

